\section{Conclusion}

A machine-learning model has no difficulty ranking one thousand stocks. An investor's problem is deciding which of those distinctions deserve capital.

This paper replaces the mechanically complete ranking with a certified partial order. Score-band separation creates a nested family of relation graphs; self-normalized calibration selects the least-abstaining graph whose future pairwise decision risk is controlled from one dependent financial history. The resulting reliable breadth measures the usable fraction of the cross-section. The edge representation maps that object directly into positions and separates certified from removed trades, while a robust pair-choice problem shows why uncertainty and trading costs generate a no-trade region.

The framework changes the economic question from ``Which model ranks stocks best?'' to ``How much of a model's ranking is reliably and implementably different?'' The empirical design is built to answer how much, when, where, and why. It evaluates whether uncertified relations are predominantly weak alpha, expensive turnover, or valuable information that the certificate removes. Each outcome is informative. The contribution is not guaranteed outperformance; it is a disciplined measurement of the gap between apparent and tradable machine-learned differentiation.
